\documentclass[11pt]{article}
\newcommand{\doctitle}{Plan: Align 2025 pm\_functions.R to Match orig (Gold Standard)}
\newcommand{\shorttitle}{Tidyverse Alignment Plan}
\newcommand{\docdate}{2026-03-24}
% pmsimstats-ng standard document preamble
% Usage: % pmsimstats-ng standard document preamble
% Usage: % pmsimstats-ng standard document preamble
% Usage: \input{pmsimstats-preamble} after \documentclass[11pt]{article}
%
% Documents must define before \input:
%   \newcommand{\doctitle}{...}       Full document title
%   \newcommand{\shorttitle}{...}     Short title for header
%   \newcommand{\docdate}{YYYY-MM-DD} Document date

\usepackage[margin=1in]{geometry}
\usepackage{amsmath,amssymb}
\usepackage[T1]{fontenc}
\usepackage{lmodern}
\usepackage{microtype}
\usepackage{graphicx}
\graphicspath{{figures/}}
\usepackage{booktabs}
\usepackage{longtable}
\usepackage{array}
\usepackage{hyperref}
\usepackage{xcolor}
\usepackage{fancyhdr}
\usepackage{parskip}
\usepackage{caption}
\usepackage{enumitem}
\usepackage{verbatim}
\usepackage{listings}

\lstset{
  language=R,
  basicstyle=\small\ttfamily,
  keywordstyle=\color{blue!70!black},
  commentstyle=\color{green!50!black},
  stringstyle=\color{red!60!black},
  breaklines=true,
  frame=single,
  framerule=0.5pt,
  rulecolor=\color{gray!40},
  backgroundcolor=\color{gray!5},
  xleftmargin=1em,
  framexleftmargin=1em,
  columns=flexible
}

\hypersetup{
  colorlinks=true,
  linkcolor=blue!60!black,
  citecolor=green!50!black,
  urlcolor=blue!60!black
}

\pagestyle{fancy}
\fancyhf{}
\lhead{\footnotesize pmsimstats-ng Technical Report}
\rhead{\footnotesize \shorttitle}
\rfoot{\footnotesize \thepage}
\renewcommand{\headrulewidth}{0.4pt}

\title{\doctitle}
\author{pmsimstats team}
\date{\docdate}
 after \documentclass[11pt]{article}
%
% Documents must define before \input:
%   \newcommand{\doctitle}{...}       Full document title
%   \newcommand{\shorttitle}{...}     Short title for header
%   \newcommand{\docdate}{YYYY-MM-DD} Document date

\usepackage[margin=1in]{geometry}
\usepackage{amsmath,amssymb}
\usepackage[T1]{fontenc}
\usepackage{lmodern}
\usepackage{microtype}
\usepackage{graphicx}
\graphicspath{{figures/}}
\usepackage{booktabs}
\usepackage{longtable}
\usepackage{array}
\usepackage{hyperref}
\usepackage{xcolor}
\usepackage{fancyhdr}
\usepackage{parskip}
\usepackage{caption}
\usepackage{enumitem}
\usepackage{verbatim}
\usepackage{listings}

\lstset{
  language=R,
  basicstyle=\small\ttfamily,
  keywordstyle=\color{blue!70!black},
  commentstyle=\color{green!50!black},
  stringstyle=\color{red!60!black},
  breaklines=true,
  frame=single,
  framerule=0.5pt,
  rulecolor=\color{gray!40},
  backgroundcolor=\color{gray!5},
  xleftmargin=1em,
  framexleftmargin=1em,
  columns=flexible
}

\hypersetup{
  colorlinks=true,
  linkcolor=blue!60!black,
  citecolor=green!50!black,
  urlcolor=blue!60!black
}

\pagestyle{fancy}
\fancyhf{}
\lhead{\footnotesize pmsimstats-ng Technical Report}
\rhead{\footnotesize \shorttitle}
\rfoot{\footnotesize \thepage}
\renewcommand{\headrulewidth}{0.4pt}

\title{\doctitle}
\author{pmsimstats team}
\date{\docdate}
 after \documentclass[11pt]{article}
%
% Documents must define before \input:
%   \newcommand{\doctitle}{...}       Full document title
%   \newcommand{\shorttitle}{...}     Short title for header
%   \newcommand{\docdate}{YYYY-MM-DD} Document date

\usepackage[margin=1in]{geometry}
\usepackage{amsmath,amssymb}
\usepackage[T1]{fontenc}
\usepackage{lmodern}
\usepackage{microtype}
\usepackage{graphicx}
\graphicspath{{figures/}}
\usepackage{booktabs}
\usepackage{longtable}
\usepackage{array}
\usepackage{hyperref}
\usepackage{xcolor}
\usepackage{fancyhdr}
\usepackage{parskip}
\usepackage{caption}
\usepackage{enumitem}
\usepackage{verbatim}
\usepackage{listings}

\lstset{
  language=R,
  basicstyle=\small\ttfamily,
  keywordstyle=\color{blue!70!black},
  commentstyle=\color{green!50!black},
  stringstyle=\color{red!60!black},
  breaklines=true,
  frame=single,
  framerule=0.5pt,
  rulecolor=\color{gray!40},
  backgroundcolor=\color{gray!5},
  xleftmargin=1em,
  framexleftmargin=1em,
  columns=flexible
}

\hypersetup{
  colorlinks=true,
  linkcolor=blue!60!black,
  citecolor=green!50!black,
  urlcolor=blue!60!black
}

\pagestyle{fancy}
\fancyhf{}
\lhead{\footnotesize pmsimstats-ng Technical Report}
\rhead{\footnotesize \shorttitle}
\rfoot{\footnotesize \thepage}
\renewcommand{\headrulewidth}{0.4pt}

\title{\doctitle}
\author{pmsimstats team}
\date{\docdate}


\begin{document}
\maketitle
\thispagestyle{fancy}

\begin{center}
\fbox{\parbox{0.95\textwidth}{\textbf{Status: Historical
planning record (2026-04-09).} The work described in this
document was largely executed in the 2026-04-08 merge that
introduced the \texttt{implementations/} directory. The
maintained tidyverse implementation now lives in
\texttt{implementations/tidyverse/R/functions.R}, and references
to \texttt{analysis/2025/pm\_functions.R} in this document
point to a path that no longer exists in the working repository
(the original 2025 pipeline has been moved to an external
archive at
\texttt{\textasciitilde/Dropbox/prj/alz/10-pmsimstats-ng/archive/2025/}).
This document is retained as a record of the planning process
and the rationale for individual alignment steps.}}
\end{center}

\bigskip
\tableofcontents

\section{Objective}

Make the 2025 tidyverse simulation pipeline produce statistically
identical results to the orig data.table pipeline for identical
parameter inputs. The 2025 codebase should be a clean, contemporary
R implementation of the same simulation -- differing only in style
(tidyverse, snake\_case, tibble), not in substance.

\section{Scope}

The target file is \texttt{analysis/2025/pm\_functions.R} (to be created
in pmsimstats-ng). The orig functions in \texttt{R/} are frozen as the
gold standard. Changes flow in one direction: 2025 adapts to orig.

\section{Naming Convention}

The 2025 code already uses snake\_case. The orig code uses camelCase
for exported functions (\texttt{buildSigma}, \texttt{generateData}, \texttt{lme\_analysis})
and internal variables. The 2025 versions will keep snake\_case
throughout. A mapping table is provided in Step 8.

\bigskip\noindent\rule{\textwidth}{0.4pt}\bigskip

\section{Step 1: Replace mod\_gompertz with modgompertz\_orig logic}

\textbf{What changes:} Replace \texttt{mod\_gompertz()} (line 12) with a
snake\_case version of the orig's \texttt{modgompertz()} that passes
through the origin.

\textbf{Current (2025):}
\begin{lstlisting}
mod_gompertz <- function(time, max_value, displacement, rate) {
  max_value * (1 - exp(-displacement * exp(-rate * time)))
}
\end{lstlisting}

\textbf{Target:}
\begin{lstlisting}
mod_gompertz <- function(t, maxr, disp, rate) {
  y <- maxr * exp(-disp * exp(-rate * t))
  vert_offset <- maxr * exp(-disp * exp(-rate * 0))
  y <- y - vert_offset
  y * (maxr / (maxr - vert_offset))
}
\end{lstlisting}

\textbf{Justification:} The current formula \texttt{max*(1 - exp(-disp*exp(-rate*t)))}
is mathematically different from the Hendrickson formulation. It does
not pass through the origin (f(0) != 0), which shifts every component
mean in the MVN distribution. This is the single largest source of
divergence between the two codebases. The \texttt{modgompertz\_orig()} that
already exists at line 1363 of pm\_functions.R proves this is known.

\textbf{Difficulty:} Trivial. Replace one function body. All callers use
the same positional arguments (time, max, disp, rate). The parameter
names differ (snake\_case vs abbreviated), but the function is called
with positional args throughout.

\textbf{Verification:} \texttt{mod\_gompertz(0, 10, 5, 0.42)} must equal 0.
\texttt{mod\_gompertz(200, 10, 5, 0.42)} must equal 10 (within tolerance).

\bigskip\noindent\rule{\textwidth}{0.4pt}\bigskip

\section{Step 2: Fix column naming in generate\_data / build\_sigma\_matrix}

\textbf{What changes:} Replace \texttt{"biomarker"} and \texttt{"baseline"} labels
with \texttt{"bm"} and \texttt{"BL"} in the sigma matrix labels, and replace
\texttt{"participant\_id"} with \texttt{"ptID"} in the output data.

\textbf{Current (2025):}
\begin{lstlisting}
labels <- c("biomarker", "baseline", ...)
\end{lstlisting}

\textbf{Target:}
\begin{lstlisting}
labels <- c("bm", "BL", ...)
\end{lstlisting}

Also in \texttt{process\_participant\_data}:
\begin{lstlisting}
participant_data <- participant_data %>%
  mutate(ptID = 1:N)
\end{lstlisting}

Also in \texttt{prepare\_analysis\_data} and \texttt{lme\_analysis}: all references
to \texttt{biomarker}, \texttt{baseline}, \texttt{symptoms}, \texttt{participant\_id} must map
to \texttt{bm}, \texttt{BL}, outcome columns, \texttt{ptID}.

\textbf{Justification:} The column names flow into the correlation matrix
row/col names, which flow into the BM-BR correlation assignment
(\texttt{correlations["bm", name1]}), which flow into the analysis model
formula, and ultimately into coefficient extraction. Mismatched
names would cause silent failures or NA results.

\textbf{Difficulty:} Moderate. Roughly 40--50 string replacements across
generate\_data, build\_sigma\_matrix, build\_correlation\_matrix,
process\_participant\_data, prepare\_analysis\_data, and lme\_analysis.
Must be done atomically -- partial changes will break the pipeline.

\textbf{Verification:} \texttt{generate\_data()} output must have columns named
\texttt{bm}, \texttt{BL}, \texttt{ptID}, \texttt{OL1}...\texttt{OL8} (matching orig).

\bigskip\noindent\rule{\textwidth}{0.4pt}\bigskip

\section{Step 3: Fix the category labels in parameter tables}

\textbf{What changes:} The 2025 resp\_param and baseline\_param tables use
long category names (\texttt{"time\_variant"}, \texttt{"pharm\_biomarker"},
\texttt{"bio\_response"}, \texttt{"biomarker"}, \texttt{"baseline"}). The orig uses short
names (\texttt{"tv"}, \texttt{"pb"}, \texttt{"br"}, \texttt{"bm"}, \texttt{"BL"}).

\textbf{Options:}

\begin{itemize}
  \item (A) Change the parameter tables passed to 2025 functions to use
    short names. This forces callers to adapt.
  \item (B) Change the 2025 functions to accept either format via a
    lookup. Adds complexity.
  \item (C) Keep long names internal to 2025 but map at boundaries.
\end{itemize}

\textbf{Recommendation:} Option A. The parameter tables are defined in
analysis scripts, not in pm\_functions.R. Change the analysis scripts
to use orig-compatible parameter tables. Then the functions can use
the same indexing as orig (\texttt{resp\_param[cat=="tv"]\$sd}).

\textbf{Justification:} Using the same parameter format eliminates an
entire class of mapping bugs and makes it possible to share
parameter definitions between the two pipelines.

\textbf{Difficulty:} Easy. The \texttt{factor\_types} / \texttt{factor\_abbreviations}
parallel arrays can be collapsed to just the abbreviations. About
20 references to long names need updating.

\textbf{Verification:} Parameter tables must be identical in structure
to those used by orig's analysis scripts.

\bigskip\noindent\rule{\textwidth}{0.4pt}\bigskip

\section{Step 4: Rewrite lme\_analysis to match orig's formula logic}

\textbf{What changes:} The 2025 \texttt{lme\_analysis()} always uses
\texttt{symptoms \textasciitilde{} biomarker + drug\_binary + biomarker:drug\_binary}.
The orig conditionally builds the formula based on design type.

\textbf{Target (matching orig logic):}

\begin{enumerate}
  \item Prepend a BL row to the trial design; compute cumulative \texttt{t}.
  \item Melt to long form; merge trial design.
  \item Compute \texttt{Dbc}:
    \begin{itemize}
      \item On-drug: \texttt{Dbc = 1}
      \item Off-drug: \texttt{Dbc = (1/2)\^{}(scalefactor * tsd / t1half)}
    \end{itemize}
  \item Test \texttt{varInDb}: does Db (on-drug indicator) vary within subjects?
  \item Test \texttt{varInExp}: does expectancy vary across timepoints?
  \item Build formula conditionally:
    \begin{itemize}
      \item Base: \texttt{Sx \textasciitilde{} bm + t}
      \item If \texttt{varInDb}: \texttt{+ Dbc + bm:Dbc}
      \item Else: \texttt{+ bm:t} (and filter to ever-on-drug participants)
      \item If \texttt{varInExp > 1} and \texttt{useDE}: \texttt{+ De}
      \item If \texttt{simplecarryover} and \texttt{varintsd}: \texttt{+ tsd}
    \end{itemize}
  \item Remove NA rows for model variables.
  \item Fit \texttt{nlme::lme} with \texttt{corCAR1}, fallback without.
  \item Extract \texttt{bm:Dbc} or \texttt{Dbc:bm} (crossover) or \texttt{bm:t}/\texttt{t:bm}
    (open label).
\end{enumerate}

\textbf{Justification:} This is the second largest source of divergence.
The formula determines which coefficient is tested for significance.
For open-label designs, orig tests a time-biomarker interaction
(\texttt{bm:t}) while 2025 tests a drug-biomarker interaction
(\texttt{biomarker:drug\_binary}). These are substantively different
hypotheses. The expectancy term \texttt{De} (actual expectancy value) vs
\texttt{t} (time) also changes model interpretation.

\textbf{Difficulty:} High. This is the most complex change. The orig
implementation is \textasciitilde 180 lines of data.table code with eval/parse
idioms for column selection. The 2025 rewrite should use
tidyverse idioms (pivot\_longer, left\_join) but replicate the
exact conditional logic. Key challenges:

\begin{itemize}
  \item The \texttt{varInDb} test requires computing per-participant mean of Db
    and checking if it is always 0 or 1. This is straightforward in
    either framework.
  \item The eval/parse column selection (\texttt{evalstring <- paste(...)}) in
    orig can be replaced with tidy selection.
  \item The \texttt{Dbc} computation must use \texttt{(1/2)\^{}(scalefactor*tsd/t1half)}
    directly, NOT lag-based propagation.
\end{itemize}

\textbf{Verification:} For identical simulated data, the two
lme\_analysis functions must return identical beta, betaSE, and
p-value (within floating-point tolerance). Test on OL, CO, and
Hybrid designs.

\bigskip\noindent\rule{\textwidth}{0.4pt}\bigskip

\section{Step 5: Match MVN draw method}

\textbf{What changes:} Replace \texttt{MASS::mvrnorm()} with Cholesky-based
draw matching orig.

\textbf{Current (2025):}
\begin{lstlisting}
participant_data <- MASS::mvrnorm(
  n = model_param$N, mu = means, Sigma = sigma,
  empirical = empirical)
\end{lstlisting}

\textbf{Target:}
\begin{lstlisting}
chol_sigma <- chol(sigma)
Z <- matrix(rnorm(n * p), nrow = n, ncol = p)
participant_data <- Z %*% chol_sigma +
  matrix(means, nrow = n, ncol = p, byrow = TRUE)
\end{lstlisting}

\textbf{Justification:} \texttt{MASS::mvrnorm} uses eigen-decomposition
internally, which consumes RNG state differently than the Cholesky
approach. For identical seeds, the two methods produce different
random draws from the same distribution. While the statistical
properties are identical in expectation, exact per-seed
reproducibility requires matching the draw method. This also
enables sigma caching (pre-compute Cholesky once, reuse across
replicates).

\textbf{Difficulty:} Easy. Replace 4 lines. The Cholesky fallback
(make.positive.definite on error) should also be copied.

\textbf{Verification:} For the same seed and sigma, both codebases
must produce identical participant data matrices.

\bigskip\noindent\rule{\textwidth}{0.4pt}\bigskip

\section{Step 6: Match PD handling in build\_sigma\_matrix}

\textbf{What changes:} The 2025 \texttt{build\_sigma\_matrix()} returns NULL
for non-PD matrices. Change to match orig: correct via
\texttt{make.positive.definite()} and continue.

\textbf{Current (2025):}
\begin{lstlisting}
if (!is_pd) {
  return(NULL)
}
\end{lstlisting}

\textbf{Target:}
\begin{lstlisting}
if (!is_pd) {
  sigma <- corpcor::make.positive.definite(sigma, tol = 1e-3)
}
\end{lstlisting}

\textbf{Justification:} Returning NULL for non-PD matrices changes the
parameter coverage of the simulation. Orig corrects and continues,
which means it can run parameter combinations that 2025 rejects.
This affects power estimates at the boundary of the PD-feasible
region (high c.bm values with OL+BDC designs).

\textbf{Difficulty:} Trivial. Replace the return(NULL) block.

\textbf{Verification:} \texttt{build\_sigma\_matrix()} with high c.bm (e.g. 0.45)
must return a valid sigma, not NULL.

\bigskip\noindent\rule{\textwidth}{0.4pt}\bigskip

\section{Step 7: Add missing infrastructure functions}

\textbf{What changes:} Port \texttt{buildtrialdesign}, \texttt{censordata}, and
\texttt{generateSimulatedResults} as tidyverse equivalents.

\subsection{7a: build\_trial\_design (from buildtrialdesign)}

The orig function (120 lines) computes tod, tsd, tpb from
timepoints, ondrug, and expectancies. The 2025 codebase has
\texttt{compute\_tod\_orig}, \texttt{compute\_tsd\_orig}, \texttt{compute\_tpb\_orig} that
replicate this logic but are not wired into a design builder.

\textbf{Approach:} Write \texttt{build\_trial\_design()} that wraps the existing
compute\_*\_orig helpers and returns a tibble-based design object
with the same structure as orig's output.

\textbf{Difficulty:} Moderate. The logic exists in pieces; needs assembly
and testing against orig's output for all four design types
(OL, CO, Hybrid, Parallel).

\subsection{7b: censor\_data (from censordata)}

The orig function (90 lines) implements MCAR+biased censoring.
Port directly with tibble/dplyr idioms.

\textbf{Difficulty:} Easy. Mostly vectorized math; the data.table
eval/parse idioms need replacement with tidy selection.

\subsection{7c: generate\_simulated\_results (from generateSimulatedResults)}

The orig function (320 lines) implements the full parameter grid
sweep with sigma caching, progressive save, and furrr parallelism.
This is the most complex function to port.

\textbf{Difficulty:} High. The nested loop structure with progressive
save, parallel dispatch, and closure-based worker functions is
intricate. However, the 2025 codebase already has furrr/future
infrastructure (unused). The port should follow orig's structure
closely.

\textbf{Justification for all three:} Without these functions, the 2025
codebase cannot run a complete simulation. It can only generate
single trials and analyze them. The simulation loop, censoring,
and trial design builder are required for the full workflow.

\bigskip\noindent\rule{\textwidth}{0.4pt}\bigskip

\section{Step 8: Style Reconciliation}

\subsection{Function name mapping}

\begin{table}[h]
\centering
\begin{tabular}{ll}
\toprule
orig (camelCase) & 2025 (snake\_case) \\
\midrule
\texttt{modgompertz} & \texttt{mod\_gompertz} \\
\texttt{cumulative} & \texttt{cumulative} (keep) \\
\texttt{buildSigma} & \texttt{build\_sigma} \\
\texttt{generateData} & \texttt{generate\_data} \\
\texttt{buildtrialdesign} & \texttt{build\_trial\_design} \\
\texttt{lme\_analysis} & \texttt{lme\_analysis} (keep) \\
\texttt{censordata} & \texttt{censor\_data} \\
\texttt{generateSimulatedResults} & \texttt{generate\_simulated\_results} \\
\texttt{validateParameterGrid} & \texttt{validate\_parameter\_grid} \\
\texttt{buildSigma} & \texttt{build\_sigma} \\
\bottomrule
\end{tabular}
\end{table}

\subsection{Variable name mapping (inside functions)}

\begin{table}[h]
\centering
\begin{tabular}{ll}
\toprule
orig & 2025 \\
\midrule
\texttt{nP} & \texttt{num\_timepoints} \\
\texttt{cl} & \texttt{factor\_abbrevs} (use \texttt{c("tv","pb","br")}) \\
\texttt{dat} & \texttt{data} \\
\texttt{modelparam} & \texttt{model\_param} \\
\texttt{respparam} & \texttt{resp\_param} \\
\texttt{blparam} & \texttt{baseline\_param} \\
\texttt{trialdesign} & \texttt{trial\_design} \\
\texttt{timeptname} / \texttt{timeptnames} & \texttt{timepoint\_name} \\
\texttt{ptID} & \texttt{ptID} (keep -- external interface) \\
\texttt{bm}, \texttt{BL} & \texttt{bm}, \texttt{BL} (keep -- external interface) \\
\bottomrule
\end{tabular}
\end{table}

\subsection{Style rules}

\begin{itemize}
  \item Use \texttt{|>} (native pipe), not \texttt{\%>\%}
  \item Use \texttt{<-} for assignment
  \item Use snake\_case for all internal variables and function names
  \item Keep camelCase for external column names that must match orig
    (\texttt{ptID}, \texttt{BL}, \texttt{Dbc}, \texttt{De})
  \item 2-space indentation
  \item Single quotes for strings
\end{itemize}

\subsection{Estimated style-only changes}

Renaming internal variables and converting pipe operators is
mechanical but pervasive. Estimate: \textasciitilde 200 line-level changes across
the full pm\_functions.R file. Low risk since these are find-replace
operations that do not affect logic.

\textbf{Difficulty:} Easy but tedious. Best done in a single pass after
all functional changes are complete.

\bigskip\noindent\rule{\textwidth}{0.4pt}\bigskip

\section{Step 9: Remove dead code from 2025}

After alignment, the following 2025-specific code can be removed:

\begin{table}[h]
\centering
\small
\begin{tabular}{lp{8cm}}
\toprule
Function / Object & Reason \\
\midrule
\texttt{mod\_gompertz} (old formula) & Replaced by origin-passing version \\
\texttt{calculate\_bio\_response\_with\_interaction} & Inline into generate\_data (orig does not have this as a separate function) \\
\texttt{calculate\_component\_halflives} & Returns uniform values; inline the single value \\
\texttt{calculate\_carryover} (multi-model) & Only exponential is used; inline the formula \\
\texttt{factor\_types} array & Collapse to \texttt{factor\_abbrevs = c("tv","pb","br")} \\
\texttt{modgompertz\_orig} & Redundant after Step 1 replaces mod\_gompertz \\
\texttt{compute\_tod\_orig} / \texttt{compute\_tsd\_orig} / \texttt{compute\_tpb\_orig} & Absorbed into build\_trial\_design \\
\texttt{build\_bm\_br\_correlations} & Absorbed into build\_correlation\_matrix \\
\texttt{build\_path\_sigma} & Redundant after primary pipeline is fixed \\
\texttt{hendrickson\_resp\_params} / \texttt{bl\_params} / \texttt{corr\_params} & Move to analysis scripts as parameter definitions \\
\texttt{validate\_correlation\_structure} & Diagnostic; keep or move to separate utils \\
\texttt{create\_sigma\_cache\_key} & Replace with orig's approach if needed \\
\texttt{report\_parameter\_validation} & Diagnostic; keep or move \\
\bottomrule
\end{tabular}
\end{table}

\bigskip\noindent\rule{\textwidth}{0.4pt}\bigskip

\section{Execution Order and Dependencies}

\begin{verbatim}
Step 1 (Gompertz)
  |
  v
Step 2 (Column names) + Step 3 (Category labels)
  |
  v
Step 5 (MVN draw) + Step 6 (PD handling)
  |
  v
Step 4 (lme_analysis)    -- depends on column names being settled
  |
  v
Step 7a (build_trial_design)
  |
  v
Step 7b (censor_data) + Step 7c (generate_simulated_results)
  |
  v
Step 8 (Style cleanup)
  |
  v
Step 9 (Dead code removal)
\end{verbatim}

Steps 1--3 can be done in a single pass (all in generate\_data
and build\_sigma\_matrix). Steps 5--6 are trivial and can be done
alongside. Step 4 is the most complex and should be done
carefully with test coverage. Steps 7a--c can proceed in parallel
once the core DGP is aligned. Steps 8--9 are cleanup.

\bigskip\noindent\rule{\textwidth}{0.4pt}\bigskip

\section{Estimated Effort}

\begin{table}[h]
\centering
\begin{tabular}{llllrr}
\toprule
Step & Description & Difficulty & Lines & Time \\
\midrule
1 & Gompertz fix & Trivial & \textasciitilde 5 & 15 min \\
2 & Column names & Moderate & \textasciitilde 50 & 1 hr \\
3 & Category labels & Easy & \textasciitilde 20 & 30 min \\
4 & lme\_analysis rewrite & High & \textasciitilde 180 & 3--4 hr \\
5 & MVN draw method & Easy & \textasciitilde 10 & 15 min \\
6 & PD handling & Trivial & \textasciitilde 5 & 10 min \\
7a & build\_trial\_design & Moderate & \textasciitilde 80 & 1.5 hr \\
7b & censor\_data & Easy & \textasciitilde 60 & 1 hr \\
7c & generate\_simulated\_results & High & \textasciitilde 250 & 4--5 hr \\
8 & Style cleanup & Easy & \textasciitilde 200 & 1.5 hr \\
9 & Dead code removal & Easy & \textasciitilde 100 & 30 min \\
\midrule
-- & \textbf{Total} & & & \textbf{\textasciitilde 14 hr} \\
\bottomrule
\end{tabular}
\end{table}

\bigskip\noindent\rule{\textwidth}{0.4pt}\bigskip

\section{Validation Strategy}

After each step, run a comparison test:

\begin{enumerate}
  \item Build identical parameter inputs (same trial design, model
    params, resp params, baseline params, seed).
  \item Call both orig and 2025 functions.
  \item Compare outputs numerically.
\end{enumerate}

\textbf{Level 1 (after Steps 1--3, 5--6):} \texttt{build\_sigma} output must
match \texttt{buildSigma} output: identical sigma matrix, means vector,
and labels.

\textbf{Level 2 (after Step 5):} \texttt{generate\_data} output must match
\texttt{generateData} output: identical participant data for same seed.

\textbf{Level 3 (after Step 4):} \texttt{lme\_analysis} output must match:
identical beta, betaSE, p for same input data, across OL, CO,
and Hybrid designs.

\textbf{Level 4 (after Steps 7a--c):} Full simulation sweep must match:
power estimates within Monte Carlo error bounds for N=200 reps.

\bigskip\noindent\rule{\textwidth}{0.4pt}\bigskip

\bigskip\noindent\rule{\textwidth}{0.4pt}\bigskip

\vfill
\noindent\rule{\textwidth}{0.4pt}
{\footnotesize
Rendered on 2026-03-25 at 19:01 PDT.\\
Source: \textasciitilde/prj/alz/10-pmsimstats-ng/pmsimstats-ng/docs/15-tidyverse-alignment-plan.tex}
\end{document}
